% Directly authored companion to gamma_structured_kernel_note.md.
\documentclass[11pt]{article}
\usepackage[T1]{fontenc}
\usepackage{amsmath,amssymb,amsthm,newtxtext,newtxmath}
\usepackage[letterpaper,margin=0.78in]{geometry}
\usepackage{graphicx,microtype,xcolor,caption,fancyhdr,tabularx}
\usepackage[colorlinks=true,linkcolor=blue!40!black,urlcolor=blue!50!black]{hyperref}
\definecolor{ink}{HTML}{263442}
\DeclareMathOperator{\sech}{sech}
\DeclareMathOperator{\csch}{csch}
\DeclareMathOperator{\sign}{sign}
\DeclareMathOperator{\diag}{diag}
\theoremstyle{definition}
\newtheorem{theorem}{Theorem}
\newcommand{\repo}{https://github.com/samlayton99/precision-mlps/blob/experiment/frozen-gamma-probe}
\newcommand{\evidence}{\repo/results/checkpoint_D_optimizers/expD36_frozen_gamma_probe/full_sweep/refinements/structured_gamma}
\graphicspath{{results/checkpoint_D_optimizers/expD36_frozen_gamma_probe/full_sweep/refinements/structured_gamma/}}
\hypersetup{pdftitle={An explicit gamma cap for finite tanh learning operators},
pdfsubject={Gamma-dependent Fourier quadratic-form bounds, target access, and small boundary-system verification}}
\setlength{\parindent}{0pt}
\setlength{\parskip}{5pt plus 1pt}
\setlength{\emergencystretch}{2em}
\setlength{\headheight}{14pt}
\captionsetup{font=small,labelfont=bf,justification=raggedright,singlelinecheck=false}
\pagestyle{fancy}\fancyhf{}
\fancyhead[L]{\small\color{ink}An explicit gamma cap for finite tanh learning}
\fancyhead[R]{\small Full proof and evidence}
\fancyfoot[C]{\small\thepage}
\renewcommand{\headrulewidth}{0.3pt}

\begin{document}
\thispagestyle{plain}
{\color{ink}\fontsize{23}{26}\selectfont\bfseries
An explicit gamma cap for\\finite tanh learning operators\par}
\vspace{8pt}
The question is whether a bound can carry gamma all the way to a learning
obstruction without hiding its effect inside a full eigendecomposition.
For a uniformly sampled, common-slope tanh dictionary, the answer is yes
at the level of quadratic forms: an explicit Fourier expression bounds
the kernel's action on every mean-zero direction. This is a finite-interval
statement about ordinary tanh features. A periodic model is not substituted
for the trained dictionary.

The distinction between this statement and a precise acquisition-time
prediction matters. The Fourier bound can be close to the actual quadratic
form while a learning-time bound using only that one number remains
conservative. We separate the analytic cap theorem, its elementary training
consequence, and a richer calculation of target energy in slow modes.
In the tested $N=512$ geometry, the cap bound is within 0.15\% of the true
quadratic form across 17 prescribed directions and four slopes. The richer
calculation gives a 5.07-million-update necessary time for the original
gamma-8 sine mixture, compared with 15.80 million updates observed.
These are different tightness statements; the latter is not an exact forecast.

\begin{center}\small
\begin{tabularx}{\linewidth}{@{}p{0.27\linewidth}X@{}}
\hline
Symbol & Meaning\\\hline
$h,q,m$ & Center spacing, integer oversampling, and observed input count\\
$\gamma,\Gamma$ & Common feature slope and its upper cap\\
$\lambda=\gamma h$ & Dimensionless bandwidth; never an eigenvalue\\
$J_\gamma,K_\gamma=J_\gamma J_\gamma^*$ & Sample-normalized features and raw-readout kernel\\
$\mu,\eta\mu$ & Kernel eigenvalue and per-update rate\\
$v=(v_i),y$ & A pattern's values at observed inputs, and normalized training target\\
$\mathcal Q_\Gamma(v)$ & Explicit cap-dependent bound on $v^*K_\gamma v$\\
$E_\gamma(n)$ & Relative residual after $n$ GD updates from zero\\\hline
\end{tabularx}
\end{center}

\section{A theorem with gamma visible in the conclusion}
Let $h>0$ and let observed inputs be a finite subset of the fine grid
$x_i=ih/q$, with positive integer $q$. There are $m$ observed inputs.
The dictionary has $W$ distinct centers from the lattice $c_j=jh$, all
with common slope $0<\gamma\le\Gamma$, and a bias. Its columns are
\[
J_\gamma[i,j]=m^{-1/2}\tanh\bigl(\gamma(x_i-c_j)\bigr),\qquad
J_\gamma[i,\mathrm{bias}]=m^{-1/2}.
\]
This includes the finite interval, endpoint sample, and uniform halo centers
in the experiments. Normalization preserves the Euclidean metric of the
original readout coefficients. With $y_i=y_i^{\rm phys}/\sqrt m$, the
relative residual from zero initialization is
$E_\gamma(n)=\|(I-\eta K_\gamma)^ny\|/\|y\|$.

A sample vector $v=(v_i)$ specifies one amplitude at each observed input
$x_i$. It can describe the training target, a residual pattern, or a chosen
oscillation. For example, given a function $f$, choose
$v_i=f(x_i)-m^{-1}\sum_r f(x_r)$ and, if desired, divide by $\|v\|$.
Subtracting the sample mean puts the vector in the theorem's domain.
Its entries are sample values, not readout weights or eigenvalues;
no kernel eigenvector is needed to choose it.

Put $v_i=0$ at unobserved fine-grid indices. Define its $q$ phase
polynomials and alias combinations, keeping dependence on $v$ explicit:
\[
V_{v,s}(\theta)=\sum_{\ell\in\mathbb Z}v_{q\ell+s}e^{-i\ell\theta},\qquad
B_{v,k}(\theta)=\sum_{s=0}^{q-1}V_{v,s}(\theta)e^{-i(\theta+2\pi k)s/q}.
\]
Combining these sums gives the exact identity
\[
B_{v,k}(\theta)=\sum_{r\in\mathrm{observed}}v_r
e^{-i(\theta+2\pi k)r/q}.
\]
This is the chosen vector's Fourier amplitude at the indicated sampled
frequency. All sample-index sums are finite. The index $k$ labels
frequencies indistinguishable when sampled at center spacing. Let
\[
M_\Gamma(\omega)=\frac z{\sinh z},\qquad
z=\frac{\pi|\omega|}{2\Gamma},\qquad M_\Gamma(0)=1,
\]
and define
\[
\mathcal Q_\Gamma(v)=\frac1{2\pi m}\int_{-\pi}^{\pi}
\left[\sum_{k\in\mathbb Z}
\frac{2M_\Gamma((\theta+2\pi k)/h)}{|\theta+2\pi k|}
\underbrace{\left|\sum_{r\in\mathrm{observed}}v_r
e^{-i(\theta+2\pi k)r/q}\right|}_{|B_{v,k}(\theta)|:\ \text{sample-vector dependence}}
\right]^2\,d\theta.
\]
This is the full oversampled formula, including all aliases; the experiment
has $q=16$. Gamma weights the Fourier content of the specified vector.
Changing $v$ changes the bound even with gamma and geometry fixed.
Scaling $v$ scales its quadratic form and bound by the squared magnitude
of that constant.

For a mean-zero vector, $B_{v,0}(0)=\sum_i v_i=0$, so the apparent
$k=0$ singularity has a finite limiting magnitude. The remaining alias
series converges exponentially. This expression uses prescribed geometry,
the target or test direction, and gamma cap. It uses no trained coefficients
or kernel eigenvectors.

\begin{theorem}[A common gamma cap limits access to resolved sample directions]
For every stated dictionary, every $0<\gamma\le\Gamma$, and every
mean-zero sample vector $v$,
\[
\boxed{v^*K_\gamma v\le\mathcal Q_\Gamma(v).}
\]
The right-hand side is nondecreasing in $\Gamma$. Its frequency weights
contain $M_\Gamma(\omega)$, which decays exponentially when
$|\omega|/\Gamma$ is large. Thus directions with little leakage into
low frequencies have a small kernel quadratic form under a small cap.

If $y\ne0$ is mean-zero, GD starts at zero, and $\eta>0$ satisfies
$\eta\|K_\gamma\|\le1$, put $r_\Gamma=\mathcal Q_\Gamma(y)/\|y\|^2$.
Whenever $\eta r_\Gamma<1$,
\[
E_\gamma(n)\ge(1-\eta r_\Gamma)^n.
\]
For $r_\Gamma>0$, reaching tolerance $0<\varepsilon<1$ requires at least
\[
\left\lceil\frac{\log(1/\varepsilon)}{-\log(1-\eta r_\Gamma)}\right\rceil
\]
updates. For $r_\Gamma=0$, a nonzero target cannot be learned.
The sufficient stable step $\eta\le1/(W+1)$ needs no spectrum calculation.
\end{theorem}
The target condition is explicit rather than an assumption about unknown
eigenvectors: sampled Fourier content enters $\mathcal Q_\Gamma$.
A high nominal oscillation frequency is insufficient; finite-window
leakage into low frequencies can dominate the integral. The theorem does
not assert that every target learns faster as gamma increases. It provides
a cap-dependent obstruction for a specified target.

\subsection*{What the directional bound says about eigenvalues}
Let $u_i$ be orthonormal eigenvectors of $K_\gamma$, with eigenvalues
$\mu_i\ge0$. For a unit mean-zero sample vector $v$,
\[
v^*K_\gamma v=\sum_i\mu_i|\langle u_i,v\rangle|^2
\le\mathcal Q_\Gamma(v).
\]
Each direction gives a different weighted average of the same eigenvalues.
A large eigenvalue is compatible with a small bound for a vector that
barely overlaps its eigenvector. The theorem controls the quadratic form
$v^*K_\gamma v$, not every component of $K_\gamma v$ or the norm
$\|K_\gamma v\|$.

If a unit eigenvector $u_j$ is itself mean-zero, choosing $v=u_j$ isolates
its eigenvalue: $\mu_j\le\mathcal Q_\Gamma(u_j)$. This requires knowing
$u_j$, which can itself change with gamma. The finite kernel need not
preserve the mean-zero subspace, so not
all its eigenvectors are admissible. Centering an eigenvector generally
destroys the eigenvector identity. The directional formula therefore does
not independently identify every eigenvalue or eigenvector.

A useful consequence needs no eigenvectors: for $t>0$,
\[
\sum_{\mu_i>t}|\langle u_i,v\rangle|^2
\le\frac{\mathcal Q_\Gamma(v)}{t}.
\]
Retaining only these nonnegative terms in the weighted average proves
the inequality. A small cap bound limits the amount of this particular
pattern that can occupy fast modes. Complementary slow energy can include
a nullspace component; Section~4 separately excludes that possibility
when bounding positive slow energy.

To guarantee many small eigenvalues, one would instead establish
$\mathcal Q_\Gamma(v)\le a\|v\|^2$ for every vector in a specified
$r$-dimensional mean-zero subspace. The min--max principle then gives
at least $r$ eigenvalues of $K_\gamma$ at most $a$, possibly including
zeros. The 17 empirical direction checks are not such a uniform subspace
proof, and their observed 0.15\% tightness is not a universal relative-error
guarantee.

\subsection*{A simpler formula at matched input and center spacing}
For $q=1$, write $V(\theta)=\sum_i v_i e^{-ii\theta}$. Pairing aliases
improves the bound to
\[
v^*K_\gamma v\le\frac1{2\pi m}\int_{-\pi}^{\pi}
\frac4{\theta^2}M_\Gamma(\theta/h)^2|V(\theta)|^2\,d\theta.
\]
The principal-frequency step envelope $4/\theta^2$ is multiplied by the
squared attenuation. The exact aligned-grid step response has squared
magnitude $\cot^2(\theta/2)$; the envelope is not that discrete spectrum.
This bounds the finite ordinary-tanh dictionary without identifying its
eigenvectors with Fourier waves.

With at least $m$ distinct tanh centers, bias and tanh columns span the
whole $m$-sample space for every $\gamma>0$. The delay theorem can then
apply to exactly representable targets. This algebraic rank statement
does not imply numerically easy interpolation or small coefficients.

\section{Proof of the cap theorem}
\textbf{Smoothing and the difference response.}
With $\rho_\gamma(x)=\frac\gamma2\sech^2(\gamma x)$, integration gives
$\rho_\gamma*\sign=\tanh(\gamma\,\cdot)$. Its Fourier transform is
$M_\gamma$. Substitute $u=e^{2\gamma x}$ in the transform integral and use
$\Gamma(1-ia)\Gamma(1+ia)=\pi a/\sinh(\pi a)$ with
$a=\omega/(2\gamma)$ to obtain the formula. The localized adjacent
feature difference
\[
d_\gamma(x)=\tanh(\gamma x)-\tanh(\gamma(x-h))
=2\rho_\gamma*\mathbf1_{[0,h]}(x)
\]
has transform
\[
\widehat d_\gamma(\omega)=
\frac{2(1-e^{-i\omega h})}{i\omega}M_\gamma(\omega).
\]
Its zero-frequency value is $2h$. Sampling this integrable difference
allows ordinary Poisson summation. Dividing out the discrete difference
factor gives, away from $\theta=0$,
\[
T_{\gamma,s}(\theta)=\sum_{k\in\mathbb Z}
\frac{2M_\gamma((\theta+2\pi k)/h)}{i(\theta+2\pi k)}
e^{i(\theta+2\pi k)s/q}.
\]
The undifferenced tanh sequence is not summable. We use this formula only
after the mean-zero sample combination cancels its constant-tail
singularity; it is not an ordinary absolutely convergent transform of
that sequence.

\textbf{Extend the centers, not the training problem.}
For mean-zero $v$, bias contributes zero. Extending center indices to all
integers adds nonnegative squared feature correlations:
\[
v^*K_\gamma v\le\frac1m\sum_{j\in\mathbb Z}
\left|\sum_i v_i\tanh\bigl(\gamma(x_i-jh)\bigr)\right|^2.
\]
The series converges: outside the sample interval, the constant $+1$ or
$-1$ parts cancel, leaving exponentially decaying correlations.
Parseval's identity expresses its right-hand side as
\[
\frac1{2\pi m}\int_{-\pi}^{\pi}
\left|\sum_{s=0}^{q-1}\overline{T_{\gamma,s}(\theta)}V_{v,s}(\theta)\right|^2d\theta.
\]
Apply the triangle inequality to the alias sum. At fixed frequency,
$M_\gamma(\omega)\le M_\Gamma(\omega)$ for $\gamma\le\Gamma$.
This proves the cap bound and monotonicity in the cap, without asserting
a Loewner ordering of the original finite kernels at different slopes.

\textbf{Matched-grid improvement.}
For $q=1$, put $a=\pi/(2\gamma h)$. Apart from the common factor
$\pi/(i\gamma h)$, the response is
\[
\sum_{k\in\mathbb Z}\csch\bigl(a(\theta+2\pi k)\bigr).
\]
For $0<\theta<\pi$, pair terms at $\theta+2\pi k$ and
$2\pi(k+1)-\theta$ to see the sum is nonnegative. Alternatively,
separate $k=0$ and pair terms at $2\pi k+\theta$ and $2\pi k-\theta$
to bound it by $\csch(a\theta)$. Odd symmetry handles negative $\theta$.
Squaring gives $4M_\gamma(\theta/h)^2/\theta^2$; replacing gamma by
its cap completes the improvement.

\textbf{Training delay.}
The target-weighted spectral measure places weight
$|\langle u_j,y\rangle|^2/\|y\|^2$ at eigenvalue $\mu_j$, including
nullspace energy. Its mean is $y^*K_\gamma y/\|y\|^2$.
The function $(1-\eta\mu)^{2n}$ is convex on $[0,1/\eta]$.
Jensen's inequality and the cap bound give
\[
E_\gamma(n)^2\ge
\left(1-\eta\frac{y^*K_\gamma y}{\|y\|^2}\right)^{2n}
\ge(1-\eta r_\Gamma)^{2n}.
\]
This proves the residual and update-count claims. Finally,
$\|K_\gamma\|\le\|J_\gamma\|_F^2\le W+1$ proves the step rule.

\textbf{Representability with enough centers.}
Set $a_i=e^{2\gamma x_i}$ and $b_j=e^{2\gamma c_j}$. Then
\[
\tanh\bigl(\gamma(x_i-c_j)\bigr)=1-\frac{2b_j}{a_i+b_j}.
\]
Bias and any $m$ tanh columns span the columns of the $m\times m$
Cauchy matrix $1/(a_i+b_j)$. Its determinant is nonzero for distinct
inputs and distinct centers. Thus the original features have full row
rank. This is not claimed for the oversampled experiment, where the
sample count exceeds the feature count.\hfill$\square$

\section{The exact Toeplitz formulation and its limits}
For consecutive uniform centers, write hidden weights as backward
differences of cumulative coefficients. Let $C$ be the invertible
transform, leaving bias unchanged, and $Z_\gamma=J_\gamma C$.
Hidden columns become the localized $d_\gamma(x-c_j)$ and one terminal
tanh feature. The metric is $T=C^*C$, with tridiagonal hidden block:
diagonal $2,\ldots,2,1$ and off-diagonal $-1$. Exactly,
\[
K_\gamma=Z_\gamma T^{-1}Z_\gamma^*,\qquad
Z_\gamma^*Z_\gamma w=\mu T w.
\]
Here $w$ is a coefficient-space generalized eigenvector. It is distinct
from the sample vector $v$ in the cap theorem: its coordinates index
transformed readout coefficients, not observed inputs.
The generalized problem recovers the same positive kernel eigenvalues.
GD in these coordinates uses $T^{-1}$ on its gradient. An identity metric
would change the optimizer. Dropping the terminal feature or restricting
to zero-sum weights also changes the problem; that coefficient subspace
is not generally invariant under raw GD.

For localized columns, summing products over the infinite input lattice
gives a Toeplitz Gram matrix. Subtracting left and right unobserved-row
Grams gives the finite matrix exactly. The outside-row feature matrices
have Hankel structure after the appropriate index reversal. Terminal
and bias cross terms remain. The discrete Toeplitz symbol sums squared
magnitudes of the $q$ sampled difference responses, divided by $m$.

This exposes gamma and supplies the cap theorem, but does not make
finite eigenvectors independent of gamma. A small absolute Toeplitz
approximation error need not be small relative to slow rates. Assess
$\widehat Z$ after metric correction:
\[
e=\|(Z_\gamma-\widehat Z)C^{-1}\|,\qquad
\|K_\gamma-\widehat K\|\le(2\|\widehat ZC^{-1}\|+e)e.
\]
The triangular factor $C^{-1}$ can amplify errors. Tests and diagnostics
retain it explicitly.

\section{Retaining target energy through small boundary systems}
The cap theorem controls a quadratic form. To retain more spectral
information, apply a smooth low-pass filter to $A=\eta K_\gamma$:
\[
f_{k,s}(a)=\frac1{1+(a/s)^k},\qquad k\text{ positive and even},\quad s>0.
\]
Define $F_{k,s}=y^*f_{k,s}(A)y/\|y\|^2$. For target energy $S(t)$ in
rates at most $t>0$, monotonicity of $f$ gives
\[
\left[\frac{F_{k,s}-f_{k,s}(t)}{1-f_{k,s}(t)}\right]_+
\le S(t)\le\min\left\{1,\frac{F_{k,s}}{f_{k,s}(t)}\right\}.
\]
Subtract an allowance on $F$ in the lower bound and add it in the upper.
These are inequalities, not an assumption that the filter is a hard
spectral projector.

Let $H=\eta\widehat J^*\widehat J$, $g=\sqrt\eta\widehat J^*y$, and
$z_j=s\exp(i(2j+1)\pi/k)$. Partial fractions and the feature-to-kernel
identity give
\[
\widehat F_{k,s}=1-\frac1{k\|y\|^2}\sum_{j=0}^{k-1}
g^*(H-z_jI)^{-1}g.
\]
The Fourier-plus-boundary construction organizes $H$ as an explicit
Fourier diagonal plus a signed low-rank correction. Woodbury's identity
reduces each shift to the retained boundary dimension. Fourier entries
and boundary factors depend explicitly on gamma; no rates are fitted
to training curves. Numerical compression and its error must be disclosed.

\textbf{Why shifted-solve residuals matter.}
In real coefficient coordinates, let $u$ approximate $(H-zI)^{-1}g$,
with $r=g-(H-zI)u$. Then
\[
g^T(H-zI)^{-1}g=g^Tu+u^Tr+r^T(H-zI)^{-1}r.
\]
The remaining error is at most $\|r\|^2/|\operatorname{Im}z|$.
These are ordinary transposes: $H$ and $g$ are real, while $u,r$ may
be complex. Transform a solve in the complex Fourier basis back before
using this identity. A small-system condition number alone is not an
accuracy certificate.

If $\|A-\widehat A\|\le\delta_A$, the resolvent identity supplies
the additional exact-arithmetic allowance
\[
|F_{k,s}-\widehat F_{k,s}|\le
\frac{\delta_A}{k}\sum_{j=0}^{k-1}
\frac{|z_j|}{|\operatorname{Im}z_j|^2}.
\]
This controls construction error separately from solve residuals. Both
sample kernels are positive semidefinite, so in this allowance replace
$|\operatorname{Im}z_j|$ by the larger distance from $z_j$ to
$[0,\infty)$. The implementation uses this improvement. Floating-point
evaluation needs a separate assessment.

\textbf{Exclude unlearnable energy without diagonalization.}
Any explicit readout $w$ gives
\[
S(0)\le\frac{\bigl(\|y-\widehat Jw\|+\epsilon_J\|w\|\bigr)^2}{\|y\|^2},
\qquad\|J_\gamma-\widehat J\|\le\epsilon_J.
\]
For $0<b<1$, put $p=[\underline S(b)-\overline S(0)]_+$ using these
bounds. Under the same stable positive step, this mass lies in $(0,b]$,
is exactly representable on the observed grid, and implies
\[
E_\gamma(n)\ge\sqrt p(1-b)^n.
\]
Consequently, when $p>\varepsilon^2$, reaching relative residual
$\varepsilon$ requires
\[
n\ge\left\lceil
\frac{\log(\sqrt p/\varepsilon)}{-\log(1-b)}
\right\rceil.
\]
The reported necessary time is the maximum of this bound over the
prescribed cutoffs. If $p\le\varepsilon^2$, that cutoff gives no
positive delay bound.
This witness-based interval starts at zero but excludes the nullspace
by subtraction. It differs from the earlier diagnostic band
$(10^{-9},b]$. Neither says the whole target is exactly representable
in an oversampled dictionary.

\section{What the experiments establish}
\begin{figure}[htbp]
\centering
\includegraphics[width=\linewidth]{structured_gamma_three_panel.png}
\caption{\textbf{An explicit gamma mechanism and its measured limits.}
A compares the cap theorem with direct ordinary-tanh quadratic forms for
a fixed Gaussian-windowed oscillation, using the training archive's
$N=512,q=16$ geometry; the displayed cap equals the actual slope.
B and C use the original sine-mixture target. B compares positive slow
energy with lower bounds from compact boundary solves, excluding
nullspace energy through explicit readout witnesses. C distinguishes
full-spectrum acquisition forecasts and actual GD first-crossing
measurements from the new necessary-time bounds. Lines connect evaluated
slopes. Sensitivity ranges compare arithmetic allowances, not confidence
intervals. No curve is a schematic; no new GPU training was run.}
\end{figure}

\subsection*{The cap bound is sharp on the tested quadratic forms}
The spectral-form audit uses $N=512$ with $q=16$ and $q=1$, and $N=128$
with $q=1$, at slopes 8, 12, 16, and 64. Its 17 directions are the
five original analytic targets after removing sample means, plus ordinary,
Gaussian-windowed, and degree-7-polynomial-complement sines at four
frequencies. This gives 204 direction--dictionary checks. Centered
versions of originally non-mean-zero targets are direction diagnostics;
they do not silently replace original training targets.

Extended-center direct sums agree with the Fourier integral to
approximately $10^{-12}$ relative error. An analytic exponential bound
controls remaining center tails below $10^{-26}$. The integral uses
8,192 and 16,384 midpoint nodes; retained aliases have a separate
analytic tail bound. These verify numerical consistency, not
outward-rounded interval certification of the integral.

For the original sine mixture at $N=512,q=16$, the infinite-center form
exceeds the finite form by factors about 1.0000608, 1.0000033,
1.0000002, and 1.0000000 as gamma increases. Across all 17 directions
on this geometry, the cap envelope is within 0.15\% of the finite form.

The clearest example is the fixed, centered, unit-norm sample vector of
\[
f(x)=e^{-\frac12(x/0.25)^2}\sin(20\pi x),\qquad -1\le x\le1.
\]
Its actual kernel quadratic form increases from $1.1323\times10^{-7}$
at gamma 8 to $0.062783$ at gamma 64, a factor over 550,000. The cap
tracks this closely. An unwindowed sine at the same frequency has
substantial finite-window low-frequency leakage; its form changes only
from about 0.243 to 0.321. The Fourier expression predicts this
distinction and shows why nominal frequency alone is inadequate.

On matched $N=512,q=1$, all sample targets are exactly representable
by the Cauchy argument. For the Gaussian-windowed target and common
step $\eta=0.5/560$, evaluating the elementary cap/Jensen formula gives
necessary times of approximately \textbf{46.0 billion}, \textbf{1.51
billion}, \textbf{51.1 million}, and \textbf{82,301} updates at caps
8, 12, 16, and 64. These are numerical lower-bound evaluations, not
executed times. They do not measure temporal tightness or prove an
ordering of actual acquisition times between these dictionaries.

\subsection*{Compact target bounds improve the delay estimate, with slack}
The primary archived target is
\[
y^{\rm phys}(x)=\sin(2\pi x)+\tfrac12\sin(6\pi x)+\tfrac14\sin(10\pi x).
\]
At $N=512,q=16$, there are 8,193 samples and 559 tanh features plus
bias. Archived steps satisfy $\eta\|K_\gamma\|\simeq0.5$ and vary with
gamma. The compact calculation uses order-64 filters, 97 logarithmic
cutoffs from $10^{-8}$ to $0.5$, scale $s=b/1.25$, and readout witnesses
from ridge shifts $10^{-10},10^{-12},10^{-14}$. It selects the largest
necessary time using only the constructed operator and these filters.
Original-kernel eigenspectra and GD crossings are afterward references.

Under the default arithmetic sensitivity allowance, necessary times at
gamma 8, 12, 16, and 64 are \textbf{5,072,048}, \textbf{130,057},
\textbf{29,640}, and \textbf{5,421}. Observed crossings are
\textbf{15,798,313}, \textbf{186,057}, \textbf{61,792}, and
\textbf{16,013}. The bounds account for about 32\%, 70\%, 48\%, and
34\% of measured time. They retain a substantial part of the delay,
while remaining conservative.

At gamma 8 the selected cutoff is $b=1.10298\times10^{-7}$. The
lower bound puts at least $0.00030614$ target energy in $(0,b]$, after
subtracting nullspace upper bound $7.25\times10^{-8}$. This is about
0.0306\% positive slow energy, rather than the larger mass at the
earlier, faster cutoff $10^{-6}$. The smaller cutoff gives the stronger
necessary time. Omitting the heuristic arithmetic allowance gives
6,978,540 updates; increasing it tenfold gives 3,805,707. These show
sensitivity, not probabilistic uncertainty.

Primary boundary solves retain dimensions 26, 26, 26, and 22. Setup
is larger: signed-correction diagonalizations have dimensions 306, 266,
246, and 170 before the small shifted solves. Numerical compression
uses a $10^{-12}$ discarded-eigenvalue threshold and is checked through
residuals. It is not a theorem that every geometry or gamma has such
a small rank. The candidate avoids a full kernel eigendecomposition,
but does eigendecompose the signed correction during setup.

Separate small-system inertia calculations bracket eigenvalues 10, 26,
and 40 in all 18 dictionaries. All 54 numerical brackets contain
independent references at each sensitivity level. At gamma 8, the 26th
eigenvalue's default relative bracket width is about 0.085\%.
The 40th is unresolved: its bracket includes zero. A reciprocal-eigenvalue
version initially gave inaccurate counts; scaling the signed correction
and stopping refinement when pivot signs are uncertain avoids falsely
narrow brackets. These checks do not certify floating-point inertia.

\subsection*{Broader checks and failed simplifications}
The sweep covers 18 dictionaries and five original targets each.
Twelve are the existing $N=128,256,512$ by gamma $8,12,16,64$ cases.
Six additions are $N=512$ at gamma $10,24,48$, and
$(N,\gamma)=(128,2),(256,4),(256,32)$. Fresh cases use the analytic
stable step $0.5/(W+1)$. Every case also reports a common-step comparison
at that rule. No target coefficients or slopes are trained in this analysis.

All 90 cases yield nonzero necessary times under the scan. Every
resolved reference crossing is later than its bound, and all 19 available
executed crossings agree with that ordering. The archived quadratic/gamma-12
trajectory remains censored at 200,000 updates. Two gamma-2 reference
crossings exceed the spectral-reference search cap $10^{18}$; these are
unresolved, not declared unlearnable. All 1,796 saved primary GD checkpoints
are retained; the full constructed-spectrum baseline matches them within
$4.56\times10^{-14}$ in relative residual.

Sharpness is not uniform. At $N=256$, gamma 4, the primary default
bound is about 302 million updates versus about 937 billion in the
numerical full-spectrum reference. The smallest allowed cutoff is
selected, indicating the scan does not resolve that slower scale.
Precision, filter coverage, and capacity witness quality matter there.
Fixed-$\gamma h$ width tracks are a limitation check as well as a scaling
check. Fixed attenuation does not imply width-independent raw rates:
feature mass and physical target frequencies also affect comparison.

Simpler candidates are unreliable substitutes. At gamma 8, a full
Fourier-diagonal model predicts 3,035 updates. A 16-mode cosine-block
approximation happens to predict 17.0 million there, but predicts 1.61
million at gamma 12 versus 186,057 measured. A fortuitous one-slope match
is insufficient. These are approximation controls, not theorem bounds.
The exact Toeplitz identity is verified separately; finite boundaries
and the original metric cannot be discarded just because the
representation is structured.

In the gamma-8 diagnostic, a Toeplitz-minus-boundary Gram discrepancy
around $7.5\times10^{-17}$ before metric correction becomes
$1.4\times10^{-12}$ afterward. Error estimates must use the raw metric.
Aliases can have large relative effects near Nyquist even if absolutely
tiny. No uniform relative alias approximation or target-independent
eigenvector stability is assumed.

The primary cases were used to develop the filter range before applying
it unchanged elsewhere. Existing results are retrospective optimization
evidence. Additional cases test transfer of the frozen calculation,
not held-out predictive generalization. Full-spectrum reference solves
are validation calculations and remain separately labeled.

\appendix
\section{Construction of the boundary systems}
This appendix specifies the finite approximation used in the compact
calculation. The cap theorem itself does not depend on it.

Let $L=Nh$. A $2L$-periodic, antiperiodic reference has absolutely
convergent expansion
\[
\phi_\gamma(d)=\sum_{k\in\mathbb Z}
\frac2{i\pi(2k+1)}M_\gamma\!\left(\frac{(2k+1)\pi}{L}\right)
e^{i(2k+1)\pi d/L}.
\]
For $|d|<L$, the difference from ordinary tanh is exactly
\[
\tanh(\gamma d)-\phi_\gamma(d)=
2\sum_{r=1}^\infty\frac{(-1)^{r-1}}{1+e^{-2r\gamma L}}
\left[e^{-2r\gamma(L-d)}-e^{-2r\gamma(L+d)}\right].
\]
To prove it, the paired image sum
\[
\tanh(\gamma d)+\sum_{\ell\ge1}(-1)^\ell
[\tanh(\gamma(d-\ell L))+\tanh(\gamma(d+\ell L))]
\]
has the derivative of $\phi_\gamma$ and the same zero value at $d=0$.
Expand its positive-distance tanh terms into convergent geometric series
to obtain the formula. Each retained term separates into an input factor
and a center factor for $d=x-c$.

Choose integer $1\le B<N/2$ and nonnegative integer $p$. Treat first
and last $B$ core-center columns exactly; retain $p$ image terms in the
remaining columns. Keep halo columns, bias, and endpoint row exact.
If $\tau_Q$ bounds the absolute sum of omitted Fourier coefficients,
\[
\epsilon_J=\sqrt{N-2B}\left[
\frac{4e^{-2(p+1)\gamma Bh}}{1-e^{-2(p+1)\gamma L}}+\tau_Q\right]
\]
bounds feature operator error by Frobenius norm. For harmonics
$k=-Q,\ldots,Q-1$ with positive integer $Q$, one valid tail bound is
\[
\tau_Q\le\frac{4\pi}{\gamma L}
\frac{e^{-a(2Q+1)}}{(1-e^{-2a(2Q+1)})(1-e^{-2a})},
\qquad a=\frac{\pi^2}{2\gamma L}.
\]
Both tails follow by bounding and summing the first omitted exponential.
Since $\|J_\gamma\|\le\sqrt{W+1}$, kernel error is at most
$(2\sqrt{W+1}+\epsilon_J)\epsilon_J$.

On the first $qN$ inputs and $N$ core centers, the reference is
diagonalized in center coordinates by the fixed half-odd Fourier matrix
$F_{\ell k}=N^{-1/2}e^{i(2k+1)\pi\ell/N}$. Its phase eigenvalues are
\[
t_{s,k}=N\sum_{a\in\mathbb Z}
\frac{2M_\gamma((2(k+aN)+1)\pi/L)}{i\pi(2(k+aN)+1)}
e^{i(2(k+aN)+1)\pi s/(qN)},
\]
and its coefficient-Gram diagonal is $m^{-1}\sum_s|t_{s,k}|^2$.
This infinite alias sum describes the exact reference. Computational $J_0$
retains harmonics $-Q,\ldots,Q-1$, so its diagonal restricts the same sum to
$-Q\le k+aN\le Q-1$; $\tau_Q$ covers omitted terms.
Sum aliases before squaring magnitudes. Zero padding includes remaining
columns and the endpoint convention.

Write the retained correction as $UV^*$, so $\widehat J=J_0+UV^*$.
Multiplication gives
\[
\widehat J^*\widehat J=J_0^*J_0+
\begin{bmatrix}J_0^*U&V\end{bmatrix}
\begin{bmatrix}0&I\\I&U^*U\end{bmatrix}
\begin{bmatrix}J_0^*U&V\end{bmatrix}^*.
\]
A QR factorization and signed eigendecomposition give the diagonal-plus-
low-rank form for Woodbury solves. Compression is a computational choice.
The target-filter residual is assessed against the uncompressed constructed
Gram matrix: a poor compressed solve enlarges its residual error instead
of silently changing the theorem's kernel.

For eigenvalue brackets, write the compressed Gram $D+VSV^*$ with $S$
diagonal with entries $\pm1$, absorbing correction magnitudes into $V$.
Away from reference poles and zero pivots, Schur-complement congruence gives
\begin{align*}
n_-(D+VSV^*-tI)
={}&n_-(D-tI)\\
&+n_-\!\left(-S-V^*(D-tI)^{-1}V\right)-n_-(-S).
\end{align*}
Here $n_-$ counts strictly negative eigenvalues. This reduces a count
query to retained boundary dimension. Compression and construction
errors must widen brackets for the actual kernel. Floating-point decisions
near poles need their own assessment; algebra alone does not certify
computed counts.

\section{Numerical status and reproduction}
The exact-arithmetic theorems, alias and image-series bounds, and
residual-squared solve identity are proved above. The numerical study
uses FP64. Its default feature sensitivity allowance is
\[
64u\sqrt{W+1}\,[1+\gamma L+\log_2N],
\]
where $u$ is machine epsilon. A shifted-solve residual-evaluation and
dot-product guard is added. These are disclosed sensitivity models,
not a complete outward-rounded analysis of FFTs, transcendental functions,
matrix products, and quadrature. The 0, 1, and 10 times feature/solve-guard
variants show their effect. None of the new target-mass bounds or timing
endpoints is labeled interval-certified. Earlier independently certified
full-spectrum timing intervals remain separate.

An independent 80-digit check recomputes selected boundary solves from
fixed compressed FP64 inputs. At the examined gamma-8 pole, the small
system's condition number is about $5.73\times10^{11}$. Residual correction
reduces its quadratic error from $1.77\times10^{-10}$ to
$6.82\times10^{-14}$. At gamma 64, ordinary rounding exceeds the pure
residual-squared radius, demonstrating the need for the separate allowance.
This checks selected solve arithmetic, not transcendental input construction
or a whole-theorem interval evaluation.

Evidence records input and code hashes, geometry, targets, steps, cutoffs,
witnesses, ranks, and numerical allowances. All 24 fingerprinted source
artifacts and earlier reports are unchanged. Independent reproduction
matches all 18 case files except elapsed time, every saved array, and both
figure files. All 25 new tests pass. The full non-slow suite reports
810 passed, 17 failed, 9 skipped, and 4 deselected; the 17 failure
identifiers exactly match the pre-existing baseline. This study used no
new GPU training. Reproduce from the repository root with one BLAS thread:
\begin{quote}\footnotesize
\begin{verbatim}
export OPENBLAS_NUM_THREADS=1
export VECLIB_MAXIMUM_THREADS=1
python -m experiments.expD36_frozen_gamma_probe.structured_symbol_audit
python -m experiments.expD36_frozen_gamma_probe.structured_gamma_analysis
python -m experiments.expD36_frozen_gamma_probe.structured_precision_audit
python -m pytest -q tests/test_expD36_structured_gamma.py \
  tests/test_expD36_structured_resolvent.py tests/test_expD36_structured_symbol.py
latexmk -pdf -interaction=nonstopmode -halt-on-error \
  -outdir=/tmp/gamma-structured-latex docs/gamma_structured_kernel_note.tex
\end{verbatim}
\end{quote}

\section*{References and relation to earlier work}
Toeplitz quadratic-form machinery is classical; see
\href{https://ee.stanford.edu/~gray/CIT006-journal.pdf}{Gray,
\emph{Toeplitz and Circulant Matrices: A Review}}.
The smoothing identity and finite Fourier boundary construction are derived
in the \href{\repo/docs/gamma_uniform_grid_theorem.md}{existing technical note}.
The new cap statement extends centers only inside an upper bound and
retains sampling aliases. Small-system filtering separately retains target
information; it does not turn the elementary Rayleigh bound into an
exact time law.
\end{document}
